% Original Markdown SHA-256: 8f389209ad2fec65f7a5498c26fc3c78cbb68a6b3bf5b8629fa22c74e43b8346
\chapter{Recurrence sources with saturated short relations: exact progression control and a seven-term obstruction}
\label{chapter:14_recurrence-control}
{\small\noindent\textbf{Source:} \nolinkurl{history/EP817_RECURRENCE_CONTROL_20260915/payloads/ep817/research/recurrence_control/notes/recurrence-control.md}\
\textbf{Identity:} \nolinkurl{8f389209ad2fec65f7a5498c26fc3c78cbb68a6b3bf5b8629fa22c74e43b8346}}
\medskip
15 September 2026. PolyClank research continuation for Erdős Problem 817.

\section{1. Scope, definitions, and outcome}\label{n14-scope-definitions-and-outcome}

For a finite positive generator set A retain its binary source, numerical image, and all representations:

\[
\begin{adjustbox}{max width=.92\linewidth}
$\displaystyle \Omega_A=\{0,1\}^{A},\qquad
\Phi_A(\epsilon)=\sum_{a\in A}\epsilon_a a,\qquad
H(A)=\Phi_A(\Omega_A).$
\end{adjustbox}
\]

A nonconstant k-term arithmetic progression is a tuple (x,x+d,...,x+(k-1)d) with d nonzero. Counts below use positive d; reversing a tuple recovers the negative-step version. The empty subset is an actual source point, with numerical value zero.

This contribution develops an arithmetic source outside the previously used fixed integer-base blocks. It proves progression exclusion at every length for a quantitative gap family, computes its exact optimal cost, and specializes to the Pell weights 1,2,5,12,29,... . For those weights it gives a complete five-state algebraic carry alphabet, an exact rational formula for four-term progression counts, and the complete integral short-relation presentation. The infinite relation operator has a sharp inverse norm 1/2. A delayed generalization has inverse norm of order the square of its delay, with explicit bounds uniform in source length and coefficient-space dimension.

An attempted acceleration of these recurrences has an unconditional obstruction: a six-generator relation pattern constructs seven actual subset sums in progression. The obstruction excludes every longer-lag recurrence of the specified form from supplying a five-, six-, or seven-term-free infinite construction, irrespective of its starting positive seeds.

No improved global numerical bound is claimed. In particular 1+sqrt(2) is larger than the preceding workbench upper rates for k=5 and k=6. The point is an all-rank test of which information an extremal lower argument must retain: an arbitrarily large admissible source can have a rank-one scalar quotient, an integrally saturated relation lattice generated by coefficient-two relations, a uniformly bounded inverse, and no five-term progression. Its actual feasible binary source and arithmetic observation cannot be recovered from those linear invariants alone.

All general statements have the proofs below. The finite graphs and calibration domains have exact executable checks and a separately implemented audit. No new Lean elaboration or global literature-priority claim is made. The prior mathematical construction attributed to the anonymous/deleted contributor and sneed-and-feed\textquotesingle s Lean extension keep their separate credits.

\section{2. A gap theorem for the complete infinite family}\label{n14-a-gap-theorem-for-the-complete-infinite-family}

Let a\_0,a\_1,... be positive integers. Put S\_j=sum\_(i=0)\^{}j a\_i and S\_j=0 for j\textless0. For a fixed integer r\textgreater=1 require

\[
\begin{adjustbox}{max width=.92\linewidth}
$\displaystyle a_j>S_{j-1}\quad(1\le j<r),\qquad
 a_j>S_{j-1}+S_{j-r}\quad(j\ge r).$
\end{adjustbox}
\tag{2.1}
\]

These inequalities retain the actual weights and all finite sums. They imply a\_j\textgreater S\_(j-1) for every j\textgreater=1, so the binary evaluation is injective: the largest differing coordinate outweighs the total contribution below it.

\textbf{Theorem R1 (uniform gap exclusion).} Every finite prefix of a sequence satisfying (2.1) has subset sums free of nonconstant progressions of length 2\^{}r+1. The conclusion covers all choices of the positive gaps and all prefix lengths simultaneously.

\textbf{Proof.} Reverse a negative-step tuple, so its step is positive. If a purported tuple has some equal highest binary coordinates, remove those common generator choices by the actual translation map. Let h be its largest differing coordinate. If h\textless r, the remaining image has at most \(2^{h+1}\leq 2^r\) values and cannot contain the tuple.

Suppose h\textgreater=r. Set T=S\_(h-r). Decompose every actual binary word into its top r coordinates, indexed h-r+1,...,h, and its remaining lower word. This gives a cover of the numerical image by 2\^{}r translates of H(\{a\_0,...,a\_(h-r)\}), each contained in an interval of width T. The word evaluation, with both parts retained, is the specified bijection onto H(\{a\_0,...,a\_h\}); its inverse reads the unique binary representation.

The tuple uses both values of the h-th binary coordinate. Its consecutive entries that cross from coordinate zero to coordinate one cross the gap between S\_(h-1) and a\_h. By (2.1), that gap is strictly larger than T. Hence its common difference exceeds T. On the other hand, 2\^{}r+1 entries in the 2\^{}r top-word classes place two entries in a single interval of width T; their positive difference is an integer multiple of the common difference, which is then at most T. This contradiction proves the theorem. Every class label and its lower source remains in the partition. QED.

For r=2 the exact condition is

\[
\begin{adjustbox}{max width=.92\linewidth}
$\displaystyle a_1>a_0,\qquad
 a_j>a_{j-1}+2S_{j-2}\quad(j\ge2).$
\end{adjustbox}
\tag{2.2}
\]

It is useful to retain the gap itself:

\[
\begin{adjustbox}{max width=.92\linewidth}
$\displaystyle g_j=a_j-a_{j-1}-2S_{j-2},\quad j\ge1.$
\end{adjustbox}
\]

Then the identity

\[
\begin{adjustbox}{max width=.92\linewidth}
$\displaystyle g_j-g_{j-1}=a_j-2a_{j-1}-a_{j-2}$
\end{adjustbox}
\tag{2.3}
\]

holds for j\textgreater=2. Thus every sequence with 0\textless a\_0=2a\_(j-1)+a\_(j-2) is in the five-term-free family. Its gaps are nondecreasing and positive. More generally the positive gaps may decrease or oscillate; only their actual positivity is used in (2.2).

\section{3. Exact extremizers inside the gap family}\label{n14-exact-extremizers-inside-the-gap-family}

Define

\[
\begin{adjustbox}{max width=.92\linewidth}
$\displaystyle p_j^{(r)}=2^j\quad(0\le j<r),\qquad
 p_j^{(r)}=2p_{j-1}^{(r)}+p_{j-r}^{(r)}\quad(j\ge r).$
\end{adjustbox}
\tag{3.1}
\]

For r=1 this is p\_j=3\^{}j. For r\textgreater=2 it is a recurrence with local coefficients of absolute value at most two.

\textbf{Theorem R2 (sharp cost for the specified gap family).} At every coordinate j, (3.1) is the smallest possible positive integral value among all sequences satisfying (2.1). For a prefix of length n, its longest progression has exactly \(\min(2^r,2^n)\) terms. Its exponential maximum-generator rate is the unique real root rho\_r\textgreater2 of

\[
\begin{adjustbox}{max width=.92\linewidth}
$\displaystyle \rho_r^r-2\rho_r^{r-1}-1=0.$
\end{adjustbox}
\tag{3.2}
\]

\textbf{Proof.} During the first r positions, the smallest positive strict superincreasing choice is 1,2,...,2\^{}(r-1). At j=r the gap p\_j-S\_(j-1)-S\_(j-r) equals one. Subtracting successive gaps gives exactly p\_j-2p\_(j-1)-p\_(j-r), so all later gaps remain one. Coordinatewise induction in (2.1), with positive integer gaps at least one, proves minimality.

The first min(r,n) generators have all subset sums from 0 through 2\^{}min(r,n)-1. Theorem R1 excludes one more term when n\textgreater=r, and cardinality excludes one more when n\textless r. This proves the exact longest length.

The function x\^{}(r-1)(x-2) is strictly increasing on x\textgreater2, starts at zero, and tends to infinity, giving the unique root. Choose positive constants c,C bounding the finitely many initial ratios p\_j/rho\_r\^{}j. The positive recurrence and its characteristic identity preserve \(c\cdot\rho_r^j\leq p_j\leq C\cdot\rho_r^j\). Taking j-th roots proves the rate. QED.

The generating function is literally

\[
\begin{adjustbox}{max width=.92\linewidth}
$\displaystyle \sum_{j\ge0}p_j^{(r)}z^j=\frac1{1-2z-z^r}.$
\end{adjustbox}
\tag{3.3}
\]

As r grows, rho\_r decreases to two; the exact equation rho\_r-2=rho\_r\^{}(1-r) gives 0\textless rho\_r-2\textless=2\^{}(1-r). This decrease comes with an increasing permitted progression length 2\^{}r. It does not improve the known large-k global upper scale, which is much smaller than a rate tending to two.

For r=2 write p\_j=p\_j\^{}(2), so

\[
\begin{adjustbox}{max width=.92\linewidth}
$\displaystyle p_0=1,\ p_1=2,\ p_{j+2}=2p_{j+1}+p_j,\qquad
 \alpha=1+\sqrt2,\quad\beta=1-\sqrt2,$
\end{adjustbox}
\]

\[
\begin{adjustbox}{max width=.92\linewidth}
$\displaystyle p_j=\frac{\alpha^{j+1}-\beta^{j+1}}{2\sqrt2}.$
\end{adjustbox}
\tag{3.4}
\]

Thus g\_5(n)\textless=p\_(n-1), while the exact minimum within (2.2) is p\_(n-1). The inclusion of the gap family into all five-admissible sources proves this upper comparison, not its reverse.

\section{4. Full integral saturation by short original relations}\label{n14-full-integral-saturation-by-short-original-relations}

For r\textgreater=2 and a length n, keep U\_n=Z\^{}(n-1), V\_n=Z\^{}n, and

\[
\begin{adjustbox}{max width=.92\linewidth}
$\displaystyle \Phi_n(z)=\sum_{j<n}p_j^{(r)}z_j.$
\end{adjustbox}
\]

Let the j-th original boundary column, 1\textless=j\textless n, be

\[
\begin{adjustbox}{max width=.92\linewidth}
$\displaystyle b_j=e_j-2e_{j-1}-\mathbf1_{j\ge r}e_{j-r},
\qquad d_n(c)=\sum_{j=1}^{n-1}c_jb_j.$
\end{adjustbox}
\tag{4.1}
\]

Each column has at most three actual coordinates and coefficients in {[}-2,2{]}.

\textbf{Theorem R3 (original scalar presentation).} For every n and r\textgreater=2,

\[
\begin{adjustbox}{max width=.92\linewidth}
$\displaystyle 0\longrightarrow U_n\xrightarrow{d_n}V_n
 \xrightarrow{\Phi_n}\mathbb Z\longrightarrow0$
\end{adjustbox}
\tag{4.2}
\]

is split exact over the integers. In particular the entire scalar kernel is generated by coefficient-two relations, the quotient is torsion-free of rank one, and these facts coexist with the all-length exclusion in Theorem R1.

\textbf{Proof and primitive.} The recurrence gives Phi\_n b\_j=0. Eliminate the largest coordinate of z using its actual column b\_j, proceeding downwards. The residual is Phi\_n(z)e\_0 because p\_0=1. The same operation records unique integral coefficients c\_j. Therefore

\[
\begin{adjustbox}{max width=.92\linewidth}
$\displaystyle z=d_n c+\Phi_n(z)e_0.$
\end{adjustbox}
\tag{4.3}
\]

The uppermost nonzero c\_j gives a nonzero j-th coordinate in d\_n c, proving injectivity and uniqueness. Equation (4.3) supplies both the scalar section and the original-boundary primitive. QED.

At r=2 this gives, for every n, a five-admissible set with

\[
\begin{adjustbox}{max width=.92\linewidth}
$\displaystyle \operatorname{span}_{\mathbb Z}
 \bigl(\ker\Phi_n\cap[-2,2]^n\bigr)=\ker\Phi_n,
 \qquad\operatorname{rank}(V_n/\ker\Phi_n)=1.$
\end{adjustbox}
\tag{4.4}
\]

There can therefore be no inequality n\textless=C(k)*rank(V\_n/ker Phi\_n) for all k-admissible sets, even when the full kernel is generated by the stated short relations. This is an explicit counterexample family to that proposed dimension route, not a counterexample to Erdős Problem 817. The feasible binary source has 2\^{}n distinct image values and remains part of the presentation.

\section{5. A uniform infinite inverse, with its delay dependence proved}\label{n14-a-uniform-infinite-inverse-with-its-delay-dependence-proved}

Let E be any real or complex Hilbert coefficient space with its specified inner product, and let S be the unilateral right shift on ell\^{}2(N\_0;E). Identify x\_j with the original c\_(j+1). For r\textgreater=2, put m=r-1. The zero-padded boundary operator is exactly

\[
\begin{adjustbox}{max width=.92\linewidth}
$\displaystyle \mathcal D_r=-2I+S-(S^*)^m.$
\end{adjustbox}
\tag{5.1}
\]

Its finite-support restriction is (4.1), including the initial columns where the shifted coordinate is absent. It is a bounded map on the entire specified Hilbert space.

\textbf{Theorem R4 (two-sided uniform control).} The operator is bijective and

\[
\begin{adjustbox}{max width=.92\linewidth}
$\displaystyle \boxed{\|\mathcal D_r^{-1}\|\le\frac{(r-1)^2}{2}.}$
\end{adjustbox}
\tag{5.2}
\]

For a nonzero coefficient space E and r=2, equality holds: the inverse norm is exactly 1/2. For nonzero E and r\textgreater=3,

\[
\begin{adjustbox}{max width=.92\linewidth}
$\displaystyle \boxed{\|\mathcal D_r^{-1}\|\ge
 \frac1{2(1-\cos(\pi/(r-2)))}
 \ge\frac{(r-2)^2}{\pi^2}.}$
\end{adjustbox}
\tag{5.3}
\]

The constants are independent of the length, the coefficient dimension, and the chosen orthonormal coordinates. The original coefficient inner product is retained. On the zero coefficient fibre the unique map is still a bijection and has norm zero; the lower bounds and norm-equality assertions explicitly require E nonzero.

\textbf{Proof.} Since S is an isometry, direct expansion gives

\[
\begin{adjustbox}{max width=.92\linewidth}
$\displaystyle -\operatorname{Re}\langle\mathcal D_r x,x\rangle
 =\frac12\left(\|(I-S)x\|^2+\|(I+S^m)x\|^2\right).$
\end{adjustbox}
\tag{5.4}
\]

Also

\[
\begin{adjustbox}{max width=.92\linewidth}
$\displaystyle I-S^m=(I+S+\cdots+S^{m-1})(I-S),$
\end{adjustbox}
\]

so \textbar\textbar(I-S\^{}m)x\textbar\textbar\textless=m\textbar\textbar(I-S)x\textbar\textbar. The parallelogram identity gives

\[
\begin{adjustbox}{max width=.92\linewidth}
$\displaystyle \|(I-S^m)x\|^2+\|(I+S^m)x\|^2=4\|x\|^2.$
\end{adjustbox}
\]

As m\textgreater=1, these identities imply that (5.4) is at least \(2\|x\|^2/m^2\). Cauchy--Schwarz gives \textbar\textbar D\_r x\textbar\textbar\textgreater=2\textbar\textbar x\textbar\textbar/m\^{}2, proving injectivity and closed range. The real quadratic form of its adjoint is identical; hence ker D\_r\^{}*=0 and the range is dense. Thus the closed range is the whole space, proving (5.2).

For r=2, S-S* is skew-adjoint and the bound is 1/2. Constant vectors truncated to longer intervals have boundary energy divided by their squared norm tending to zero, so \textbar\textbar D\_2 x\textbar\textbar/\textbar\textbar x\textbar\textbar{} tends to two; the inverse bound is sharp.

For r\textgreater=3 choose theta=pi/(m-1) and truncate the sequence \(\exp(i\cdot j\cdot\theta)v\) for a fixed unit coefficient v. Away from the finitely many cutoff rows, D\_r acts by

\[
\begin{adjustbox}{max width=.92\linewidth}
$\displaystyle -2+e^{-i\theta}-e^{im\theta}
 =-2+2\cos\theta,$
\end{adjustbox}
\]

because m theta=pi+theta. The number and magnitudes of boundary errors are fixed as the interval length grows. Their squared norm divided by the interval length tends to zero. Consequently the ratio \textbar\textbar D\_r x\textbar\textbar/\textbar\textbar x\textbar\textbar{} tends to 2(1-cos theta), proving the first lower bound. The inequality 1-cos t\textless=t\^{}2/2 proves the second. In a real coefficient space the same bound follows from the real and imaginary parts, or from the equality of an operator\textquotesingle s real and complexified norms. QED.

The inverse also has a convergent operator integral with an explicit uniform remainder:

\[
\begin{adjustbox}{max width=.92\linewidth}
$\displaystyle \boxed{D_r^{-1}=-\int_0^\infty e^{tD_r}\,dt,\qquad
 \left\|D_r^{-1}+\int_0^T e^{tD_r}\,dt\right\|
 \le\frac{(r-1)^2}{2}e^{-2T/(r-1)^2}.}$
\end{adjustbox}
\tag{5.4a}
\]

Indeed, differentiating \textbar\textbar exp(tD\_r)x\textbar\textbar\^{}2 and applying (5.4) gives \textbar\textbar exp(tD\_r)\textbar\textbar\textless=exp(-2t/(r-1)\^{}2). The integral converges in operator norm. Integrating its derivative proves both inverse laws, and integrating the exponential bound gives the displayed tail. This is a uniform calculation on the entire specified infinite source, not an interchange of a fixed-length calibration with an unproved limiting claim.

For r=2 and a finite primitive c\_1,...,c\_(n-1), put c\_0=c\_n=c\_(n+1)=0. The exact finite identity is particularly simple:

\[
\begin{adjustbox}{max width=.92\linewidth}
$\displaystyle \boxed{\|d_n c\|^2=4\sum_{j=1}^{n-1}\|c_j\|^2
 +\sum_{j=0}^{n-1}\|c_{j+2}-c_j\|^2.}$
\end{adjustbox}
\tag{5.5}
\]

Every boundary term is present. In matrix form the relation Gram has diagonal 5,6,...,6, zero first off-diagonals and minus one second off-diagonals. The identity includes n=1 and n=2 with their literal zero-dimensional or one-dimensional source.

\subsection{5.1 The quotient-completion comparison}\label{n14-the-quotient-completion-comparison}

On finite coefficient supports the arithmetic quotient in (4.2) is E. The Hilbert completion of (5.1) is onto and has zero quotient. The actual comparison map of complexes therefore induces E-\textgreater0. It does not say that the original supported arithmetic class was absent.

The lost observation is explicit. For fixed v in E,

\[
\begin{adjustbox}{max width=.92\linewidth}
$\displaystyle z^{(N)}=(p_N^{(r)})^{-1}e_Nv\longrightarrow0
 \quad\text{in the raw Hilbert norm},
 \qquad\Phi z^{(N)}=v.$
\end{adjustbox}
\]

Retain it through the graph source z-\textgreater(z,Phi z). Its closure is ell\^{}2(E) direct-sum E: truncate any first component, then correct its amplitude at a sufficiently high coordinate using the displayed vector. The completed relation map and contraction are

\[
\begin{adjustbox}{max width=.92\linewidth}
$\displaystyle \widehat d x=(D_rx,0),\quad\pi(z,v)=v,\quad s(v)=(0,v),\quad
 K(z,v)=D_r^{-1}z,$
\end{adjustbox}
\]

\[
\begin{adjustbox}{max width=.92\linewidth}
$\displaystyle K\widehat d=I,\qquad\widehat dK+s\pi=I.$
\end{adjustbox}
\tag{5.6}
\]

The forgetful map to the raw Hilbert completion has kernel 0 direct-sum E. These are maps on the original completed relation and observation spaces.

At finite length retain \(W_n=\sum_{j<n}(p_j^{(r)})^2\). The least-norm scalar section is a-\textgreater(p\_j a/W\_n)\_j. The raw and graph quotient Grams are

\[
\begin{adjustbox}{max width=.92\linewidth}
$\displaystyle G_n^{\rm raw}=G/W_n,\qquad
 G_n^{\rm graph}=(1+1/W_n)G.$
\end{adjustbox}
\tag{5.7}
\]

For a coefficient space of finite dimension h, the original return cost is

\[
\begin{adjustbox}{max width=.92\linewidth}
$\displaystyle \log\det G_n^{\rm raw}
 =\log\det G_n^{\rm graph}-h\log(1+W_n).$
\end{adjustbox}
\tag{5.8}
\]

For Pell weights, 2W\_n=p\_(n-1)p\_n, proved by telescoping p\_jp\_(j+1)-p\_(j-1)p\_j=2p\_j\^{}2 (and the initial endpoint). Thus the raw quotient Gram decays at the explicit exponential rate alpha\^{}(-2n), while the infinite inverse norm remains 1/2. The source\textquotesingle s uniform inverse and the arithmetic observation\textquotesingle s return cost are simultaneously present.

\section{6. The algebraic observation and its supported zero directions}\label{n14-the-algebraic-observation-and-its-supported-zero-directions}

For the Pell case let E\_Z=Z{[}alpha{]}, alpha\^{}2=2alpha+1. Use the literal basis (1,alpha) and the map

\[
\begin{adjustbox}{max width=.92\linewidth}
$\displaystyle J_n:Z^n\to E_Z,\qquad J_ne_j=\alpha^j,\qquad
 \nu(u+v\alpha)=u+2v.$
\end{adjustbox}
\tag{6.1}
\]

By (3.4), nu(alpha\^{}j)=p\_j, so Phi\_n=nu J\_n. For n\textgreater=2, J\_n is onto. Its kernel is generated by the shifted polynomial relations

\[
\begin{adjustbox}{max width=.92\linewidth}
$\displaystyle e_{j+2}-2e_{j+1}-e_j,\qquad0\le j<n-2.$
\end{adjustbox}
\]

Monic polynomial division proves this over Z, with the actual quotient coefficients retained. The additional scalar relation b\_1=e\_1-2e\_0 maps to alpha-2. Hence

\[
\begin{adjustbox}{max width=.92\linewidth}
$\displaystyle \boxed{0\to\ker J_n\to\ker\Phi_n
 \xrightarrow{J_n}\mathbb Z(\alpha-2)\to0.}$
\end{adjustbox}
\tag{6.2}
\]

The scalar observation is Z-linear, not a ring homomorphism. The element alpha-2 is a nonzero element of E\_Z with supported scalar image zero. Replacing the terminal condition nu(z)=0 by z=0 discards actual arithmetic tuples.

For example the binary representations of 0,1,2,3 in the weights 1,2 have algebraic second differences alpha-2 and 2-alpha. Both are nonzero and both have scalar image zero. This example is retained as a negative control in the finite certification.

\section{7. A complete five-state carry computation for every word length}\label{n14-a-complete-five-state-carry-computation-for-every-word-length}

For a coefficient word c\_j in {[}-2,2{]}, evaluate its polynomial in alpha from the highest coefficient to the lowest by Horner\textquotesingle s rule. In integer coordinates the transition is

\[
\begin{adjustbox}{max width=.92\linewidth}
$\displaystyle (u,v)\longmapsto(v+c,u+2v).$
\end{adjustbox}
\tag{7.1}
\]

The actual terminal arithmetic condition is u+2v=0.

\textbf{Lemma R5 (exact bounded carry alphabet).} Every Horner prefix of a scalar-zero word belongs to the following five-element set:

\[
\begin{adjustbox}{max width=.92\linewidth}
$\displaystyle \boxed{\mathcal C=\{0,1,-1,\alpha-2,2-\alpha\}.}$
\end{adjustbox}
\tag{7.2}
\]

\textbf{Proof.} Conjugation sends alpha to beta=1-sqrt(2), with \textbar beta\textbar=sqrt(2)-1. Every finite Horner prefix z obeys

\[
\begin{adjustbox}{max width=.92\linewidth}
$\displaystyle |z'|<\frac2{1-|\beta|}=2+\sqrt2.$
\end{adjustbox}
\]

If the full word z has nu(z)=0, then z=v(alpha-2), and its conjugate is -alpha*v. Therefore \textbar v\textbar\textless sqrt(2), so v is -1,0, or 1 and \textbar z\textbar\textless=alpha-2.

If a prefix has t remaining digits, the full-word equation is z=alpha\^{}t q+R, with \textbar R\textbar\textless=2(alpha\^{}t-1)/(alpha-1). Hence

\[
\begin{adjustbox}{max width=.92\linewidth}
$\displaystyle |q|\le\frac2{\alpha-1}
 +\frac{|z|-2/(\alpha-1)}{\alpha^t}<\sqrt2.$
\end{adjustbox}
\]

Write q=(u+v)+v sqrt(2), q\textquotesingle=(u+v)-v sqrt(2). The two strict bounds give \textbar v\textbar\textless1+sqrt(2), and more sharply \textbar v\textbar\textless(2+2sqrt(2))/(2sqrt(2))\textless2, so v is -1,0,1. For v=0 the possible u are -1,0,1. For v=1, intersecting the two bounds gives u=-2; the other case gives u=2. This is exactly (7.2). All strict endpoints matter. QED.

For k rows a state is (z\_0,...,z\_(k-3),flag), with z\_i in C. A binary column epsilon in \{0,1\}\^{}k updates the carries by (7.1) using c\_i=epsilon\_i-2epsilon\_(i+1)+epsilon\_(i+2). The flag records whether the first two binary words have ever differed. A target is admitted only when every updated carry belongs to C. Initial state is all zero with false flag. Acceptance requires true flag and nu(z\_i)=0 for every row.

The binary numerical evaluation is injective by Section 2. Consequently the flag is exactly numerical nonconstancy for a terminal arithmetic tuple. The path-to-tuple map evaluates the actual columns in the reversed Pell weights. The inverse reads each unique binary representation and its Horner prefixes. Lemma R5 proves that every accepted arithmetic tuple stays in the finite carrier; no real progression is removed by that bound.

The complete reachable data are:

\begin{longtable}[]{@{}rrrr@{}}
\toprule\noalign{}
k & states & labeled transitions & accepting states \\
\midrule\noalign{}
\endhead
\bottomrule\noalign{}
\endlastfoot
3 & 8 & 26 & 3 \\
4 & 22 & 46 & 3 \\
5 & 27 & 28 & 0 \\
6 & 49 & 50 & 0 \\
\end{longtable}

Every listed transition is independently re-enumerated from actual binary columns. The 27-state five-term exclusion is a second proof of the Pell specialization of Theorem R1. The finite graph, together with the prefix lemma and the exact terminal map, proves every length. It is not an extrapolation from tested prefixes.

\subsection{7.1 Exact four-term progression counts}\label{n14-exact-four-term-progression-counts}

Let F\_n be the number of positive-step four-term progressions among the subset sums of the first n Pell weights, with F\_0=0. The weighted path count gives

\[
\begin{adjustbox}{max width=.92\linewidth}
$\displaystyle \boxed{\sum_{n\ge0}F_nz^n
 =\frac{z^2}{(1-2z)(1-z-4z^2)}.}$
\end{adjustbox}
\tag{7.3}
\]

Equivalently, F\_0=F\_1=0, F\_2=1, and

\[
\begin{adjustbox}{max width=.92\linewidth}
$\displaystyle F_n=3F_{n-1}+2F_{n-2}-8F_{n-3}\quad(n\ge3).$
\end{adjustbox}
\tag{7.4}
\]

The first values are 0,0,1,3,11,31,91,247,675,1791,4747,12423,32435. The dominant counting rate is (1+sqrt(17))/2, strictly larger than the image-cardinality rate two. Thus the infinite family is richly four-progressive while remaining five-progressive-free.

For a small transparent certificate, retain the eight states that can reach acceptance. In order they are

\[
\begin{adjustbox}{max width=.92\linewidth}
$\displaystyle (r,-r;1),\ (-1,1;0),\ (-1,1;1),\ (0,0;0),\
 (0,0;1),\ (1,-1;0),\ (1,-1;1),\ (-r,r;1),
 \quad r=\alpha-2.$
\end{adjustbox}
\]

Their labeled-edge multiplicity matrix is

\[
\begin{adjustbox}{max width=.92\linewidth}
$\displaystyle T=\begin{pmatrix}
0&0&1&0&1&0&2&0\\
0&0&0&0&0&0&0&1\\
0&0&0&0&0&0&0&1\\
0&1&0&2&0&1&0&0\\
0&0&1&0&2&0&1&0\\
1&0&0&0&0&0&0&0\\
1&0&0&0&0&0&0&0\\
0&0&2&0&1&0&1&0
\end{pmatrix}.$
\end{adjustbox}
\]

The initial row is index 3, counting from zero; the accepting indices are 0,4,7. Thus F\_n=(e\_3\^{}T T\^{}n (e\_0+e\_4+e\_7))/2. Reversal pairs the two step signs. Solving (I-zT)f=e\_0+e\_4+e\_7 gives (7.3). The executable also certifies the denominator on the full 22-state graph: its recurrence-residual row has zero outputs for the first 22 powers; Cayley--Hamilton then proves all subsequent outputs vanish. The full discarded states remain in the receipt and are excluded from this smaller counting matrix only by their proved inability to reach an accepting state.

Pisot-base finite-state arithmetic is established prior methodology (Frougny \& Pelantová, 2018). Here the complete small alphabet, actual terminal observation, all labeled transitions and integer arithmetic checks are supplied for this specific recurrence. No generic literature theorem is substituted for these calculations.

\subsection{7.2 An explicit scalar-observation carrier for every delay}\label{n14-an-explicit-scalar-observation-carrier-for-every-delay}

The finite arithmetic control extends to every r\textgreater=2, not only the quadratic case. The bounds below deliberately retain their degree dependence rather than declaring a uniform compact coefficient set as r increases.

Put

\[
\begin{adjustbox}{max width=.92\linewidth}
$\displaystyle Q_r(X)=X^r-2X^{r-1}-1,\quad E_r=\mathbb Z[X]/(Q_r),\quad
 \nu_r\left(\sum_{i<r}z_iX^i\right)=\sum_{i<r}2^iz_i.$
\end{adjustbox}
\tag{7.5}
\]

This is the actual scalar observation satisfying \(\nu_r(X^j)=p_j^{(r)}\) for every j. It is Z-linear, and its kernel is retained. Let M\textgreater=1 be an integer coefficient bound, and define the explicit positive rational numbers and integer cutoff

\[
\begin{adjustbox}{max width=.92\linewidth}
$\displaystyle \delta_r=\frac{2}{(r-1)^2(3r-2)},\qquad
 B_r(M)=\frac{M(r-1)(r+2)}{2r\delta_r^2},$
\end{adjustbox}
\] \[
\begin{adjustbox}{max width=.92\linewidth}
$\displaystyle \boxed{C_r(M)=\left\lceil
 \frac{M(r-1)(r+2)2^{r-1}3^{r-2}}{r\delta_r^3}
 \right\rceil.}$
\end{adjustbox}
\tag{7.6}
\]

\textbf{Theorem R5u (all-delay finite arithmetic carrier).} Every Horner prefix of every finite coefficient word in {[}-M,M{]} whose full scalar observation nu\_r is zero has all of its original coefficients in {[}-C\_r(M),C\_r(M){]}. Consequently, for fixed r and k, all progression equations on every recurrence prefix are decidable in a finite graph with at most

\[
\begin{adjustbox}{max width=.92\linewidth}
$\displaystyle 2(2C_r(2)+1)^{r(k-2)}$
\end{adjustbox}
\]

states. Actual terminal membership in ker nu\_r is used, rather than zero in the larger algebra. This coarse bound is a proved finite domain, not a claim that its whole graph was enumerated in the delivered computation.

\textbf{Proof of the root and observation estimates.} Q\_r has exactly one root alpha\textgreater1 outside the unit circle, with 2\textless alpha\textless3, and its other r-1 roots lie strictly inside. To verify the count, for 0\textless t\textless1 apply Rouche\textquotesingle s theorem on the unit circle to \(X^r-2X^{r-1}-t\): the term -2X\^{}(r-1) has magnitude 2, greater than 1+t. Thus r-1 roots lie inside. At t=1 there is no unit-circle root: equality in \(2|z|^{r-1}=|z^r-1|\) at \textbar z\textbar=1 would force both z\^{}r=-1 and z\^{}(r-1)=-1, hence z=1, a contradiction. The root count is unchanged on reaching t=1. The positive root is unique by (3.2). The roots are simple because a common nonzero root of Q\_r and Q\_r\textquotesingle{} would be 2(r-1)/r\textless2, where Q\_r is negative. In particular alpha is a Pisot number. Irreducibility also follows: a nonconstant monic integer factor excluding alpha would have nonzero integral constant coefficient of modulus less than one, since all its roots are inside the unit circle.

On \textbar z\textbar=1, with m=r-1,

\[
\begin{adjustbox}{max width=.92\linewidth}
$\displaystyle \left|Q_r(z)\right|\ge 2-\operatorname{Re}z+\operatorname{Re}z^m
 =\tfrac12(|1-z|^2+|1+z^m|^2)\ge2/m^2.$
\end{adjustbox}
\]

The final inequality is the scalar version of (5.4): use \textbar1-z\^{}m\textbar\textless=m\textbar1-z\textbar{} and \(|1-z^m|^2+|1+z^m|^2=4\). Inside the unit disk \textbar Q\_r\textquotesingle\textbar\textless=3r-2. Applying the derivative bound along a radius from a stable root lambda to the unit circle proves

\[
\begin{adjustbox}{max width=.92\linewidth}
$\displaystyle 1-|\lambda|\ge\delta_r.$
\end{adjustbox}
\tag{7.7}
\]

The constant coefficient of Q\_r has modulus one, alpha\textless3, and all other roots have modulus less than one. Hence each stable root has modulus greater than 1/3.

Let lambda\_0=alpha and lambda\_1,...,lambda\_(r-1) be the other roots. Partial fractions of the literal generating function (3.3) give the scalar observation formula

\[
\begin{adjustbox}{max width=.92\linewidth}
$\displaystyle \nu_r(q)=\sum_{a=0}^{r-1}w_a q(\lambda_a),\qquad
 w_a=\frac{\lambda_a^{r-1}}{Q_r'(\lambda_a)}
     =\frac{\lambda_a}{r\lambda_a-2(r-1)}.$
\end{adjustbox}
\tag{7.8}
\]

Indeed expansion at zero gives this equality first on X\^{}j, hence on every original remainder. The expanding coefficient satisfies w\_0\textgreater=2/(r+2)\textgreater0. For the stable coefficients,

\[
\begin{adjustbox}{max width=.92\linewidth}
$\displaystyle |w_a|\le1/(r\delta_r),$
\end{adjustbox}
\]

because \textbar r lambda\_a-2(r-1)\textbar\textgreater=r-2+r delta\_r\textgreater=r delta\_r.

Every finite Horner prefix has \textbar q(lambda\_a)\textbar\textless=M/delta\_r at a stable root. If its full word z has nu\_r(z)=0, equation (7.8) therefore gives \textbar z(alpha)\textbar\textless=B\_r(M). A prefix q with t remaining digits satisfies z(alpha)=alpha\^{}t q(alpha)+R with \textbar R\textbar\textless=M(alpha\^{}t-1)/(alpha-1). Since B\_r(M)\textgreater=M and alpha\textgreater2,

\[
\begin{adjustbox}{max width=.92\linewidth}
$\displaystyle |q(\alpha)|\le B_r(M),\qquad
 |q(\lambda_a)|\le M/\delta_r\le B_r(M).$
\end{adjustbox}
\tag{7.9}
\]

Finally retain the exact interpolation map back to the original monomial coordinates:

\[
\begin{adjustbox}{max width=.92\linewidth}
$\displaystyle q(X)=\sum_{a=0}^{r-1}q(\lambda_a)
 \frac{Q_r(X)}{(X-\lambda_a)Q_r'(\lambda_a)}.$
\end{adjustbox}
\tag{7.10}
\]

Every coefficient of Q\_r(X)/(X-lambda\_a) has magnitude at most 2\^{}r, since the product of 1+\textbar lambda\_b\textbar{} over the other roots is at most \(4\cdot2^{r-2}\). At stable roots, \(|Q_r'(\lambda_a)|\geq3^{-(r-2)}\cdot r\delta_r\). The same common reciprocal bound 3\^{}(r-2)/(r delta\_r) bounds the expanding derivative reciprocal as well. Equations (7.9)-(7.10) yield the coefficient bound

\[
\begin{adjustbox}{max width=.92\linewidth}
$\displaystyle |[X^i]q|\le
 B_r(M)\,2^r3^{r-2}/\delta_r\le C_r(M).$
\end{adjustbox}
\]

This proves the finite carrier in the literal Z-basis; the complex embeddings serve only as specified estimates and are not replacements for that basis. QED.

\textbf{Proof of the progression decision statement.} Use Horner multiplication by X in E\_r, reducing by the original relation \(X^r=2X^{r-1}+1\), and add each second-difference coefficient from the actual binary k-column. Retain only target coefficient vectors in the proved box, plus the first-two-word difference flag. Every genuine progression stays in the box by the theorem; every accepted path evaluates to its actual scalar progression because nu\_r(X\^{}j)=p\_j. The binary evaluation is injective by Theorem R1, so the flag is the exact nonconstancy test. The graph is finite with the stated size bound, and reachability is equivalent to existence at some finite length. All support labels, algebraic carries, and original scalar-zero terminal classes remain explicit.

For r=2, M=2 the coarse formula gives C\_2(2)=64. Lemma R5 and its two-embedding rectangle refine that entire coarse lattice box to the exact five values used in the fully enumerated graphs. For r\textgreater2 the delivered finite checks test original polynomial reduction and scalar-zero prefix containment in bounded word domains; the complete higher-degree graph is not claimed to have been constructed. The theorem gives its explicit stopping domain.

\section{8. The faster-delay obstruction: seven actual subset sums}\label{n14-the-faster-delay-obstruction-seven-actual-subset-sums}

Consider the six distinct positive generators

\[
\begin{adjustbox}{max width=.92\linewidth}
$\displaystyle x,\ b,\ u=x+2b,\ c,\ v=b+2c,\ w=u+2v=x+4b+4c.$
\end{adjustbox}
\tag{8.1}
\]

\textbf{Theorem R6 (seven-term relation motif).} Their subset sums contain a seven-term progression of positive difference b+c. In the displayed generator order, its binary masks are

\[
\begin{adjustbox}{max width=.92\linewidth}
$\displaystyle 1,\ 11,\ 19,\ 28,\ 32,\ 42,\ 50.$
\end{adjustbox}
\tag{8.2}
\]

The actual values are

\[
\begin{adjustbox}{max width=.92\linewidth}
$\displaystyle x,\quad x+b+c,\quad x+b+v,\quad u+c+v,\quad
 w,\quad w+b+c,\quad w+b+v,$
\end{adjustbox}
\]

which equal x+j(b+c), 0\textless=j\textless=6. This proves the assertion by direct substitution. The second differences of the masks are integer combinations of the original three relation rows

\[
\begin{adjustbox}{max width=.92\linewidth}
$\displaystyle (-1,-2,1,0,0,0),\quad(0,-1,0,-2,1,0),\quad(0,0,-1,0,-2,1).$
\end{adjustbox}
\]

Those coefficients are included in the finite proof data.

Now take any strictly increasing positive initial seeds a\_0,...,a\_(r-1), r\textgreater=3, and extend by

\[
\begin{adjustbox}{max width=.92\linewidth}
$\displaystyle a_j=2a_{j-1}+a_{j-r}\quad(j\ge r).$
\end{adjustbox}
\]

At the six distinct indices

\[
\begin{adjustbox}{max width=.92\linewidth}
$\displaystyle 0,\ r-1,\ r,\ 2r-2,\ 2r-1,\ 2r$
\end{adjustbox}
\]

the weights have exactly the form (8.1). Thus the first 2r+1 generators always contain the displayed seven-term progression. This is a theorem for every choice of positive increasing seeds, not a bound on a finite seed search. In particular these longer delays cannot yield infinite k-admissible constructions for k=5,6,7, even though their characteristic rates approach two.

The index distinctness is indispensable. At r=2 the positions r and 2r-2 coincide. The original six-label map to the five actual generator coordinates is

\[
\begin{adjustbox}{max width=.92\linewidth}
$\displaystyle (0,1,2,3,4,5)\longmapsto(0,1,2,2,3,4).$
\end{adjustbox}
\tag{8.3}
\]

It is a linear map on free coefficient modules, but its pushforward sends the fourth mask in (8.2) to a vector with coefficient two at the third generator. Thus it does not send the whole binary feasible source into the five-generator binary source. For the actual Pell weights 1,2,5,12,29, the candidate seven values are 1,8,15,22,29,36,43, and 22 is absent. It lies in the retained gap (20,29). All other displayed values are present. The same scalar relation equations therefore yield an admissible witness for r\textgreater=3 but an unsupported binary request at r=2; the exact support-changing map records the distinction.

\subsection{8.1 Equal-difference packing for arbitrary progression length}\label{n14-equal-difference-packing-for-arbitrary-progression-length}

The motif also gives an all-k constraint with its support hypothesis retained. Suppose s copies of (8.1) have pairwise disjoint actual generator sets and the same positive difference D=b+c. Let their initial values be x\_1,...,x\_s. For every integer h with 0\textless=h\textless=6s, choose actual integers h\_j in {[}0,6{]} with sum h (fill the first coordinate up to six, then the next, and so on). In motif j take the subset with mask (8.2) at index h\_j. Their disjoint union is an actual generator subset of value

\[
\begin{adjustbox}{max width=.92\linewidth}
$\displaystyle \sum_{j=1}^{s}x_j+hD.$
\end{adjustbox}
\]

Thus s such motifs force a (6s+1)-term progression. A k-admissible set can contain at most floor((k-2)/6) pairwise disjoint copies with a fixed common D. This is a support-sensitive packing condition, not a bound on the number of overlapping symbolic relations. The complete incidence-labelled witnesses make the restriction usable in a candidate search or in a later matching argument.

\section{9. Consequences for the general research programme}\label{n14-consequences-for-the-general-research-programme}

Theorem R3 excludes a tempting all-rank strategy: the number of generators cannot be bounded using only the rank of the saturated short-relation quotient, even with primitive coefficient-two relations. Theorem R4 strengthens that test: a uniform source inverse does not add the missing information, since its norm is exactly 1/2 in a five-admissible family of arbitrary length. The true binary source and its observation remain 2\^{}n-point arithmetic data, and (5.8) gives the quantitative cost lost by forgetting the original observation in Hilbert completion.

The gap theorem supplies a working replacement in this family: retain the original high-coordinate partition and its actual separation budget. It handles arbitrary positive gap sequences and proves the exact best recurrence within that class. The finite algebraic state calculation supplies another replacement, capable of retaining supported scalar-zero directions such as alpha-2 rather than discarding them at a stronger observation.

The seven-term motif then tests an attempted escape from the gap optimum. The faster delayed recurrence has an explicit all-seed obstruction after at most 2r+1 generators. Any new recurrence intended to improve the k=5 or k=6 bounds must avoid this motif with its six distinct original supports, or supply a different admissibility mechanism. This is a necessary arithmetic constraint on candidate constructions; it is not asserted sufficient for general admissibility.

The unrestricted rates and their sharp large-k lower law remain unevaluated by this contribution. The inspected general benchmark gives a lower logarithmic scale of order 1/k and an upper scale of order log(k)/k (Korsky, 2026). The exact maps above identify concrete information that a proposed uniform lower argument must use, and a concrete relation motif it must recognize.

\section{10. Verification and source record}\label{n14-verification-and-source-record}

Run the standard-library producer and separately implemented auditor from this directory\textquotesingle s certificates folder, supplying the paths shown in its README. No software download, floating-point comparison, or external service is needed to check the finite data.

The producer constructs all reachable graphs from the exact five-element carry alphabet, retains every binary edge label, and checks its counts against independent integer-bitset progression counts through thirteen weights. The algebraic carry alphabet is verified by exact signs in Q(sqrt(2)); its finiteness proof is in Lemma R5. The finite checks also cover all 2,186 specified positive-gap sequences through length seven, the stated dyadic-gap prefixes, 3,240 triangular basis identities through length eighty, exact coercivity and section metrics, 90 delayed-recurrence applications through delay 32, 690 finite general-delay operator identities, and complete scalar-zero coefficient-word checks for delays two through six and word lengths zero through six. These are bounded calibrations; they are not labelled proofs by exhaustion of infinite parameter families.

The auditor imports no producer module. It reconstructs the complete carry transitions and their accepted scalar observation, compares arithmetic AP counts by direct pairs, checks original basis primitives and finite Grams independently, and verifies the seven-term masks and their alias failure. The completed audit is recorded in certificates/independent\_audit.json. Normal/optimized comparisons, mutation controls, and precise evidence boundaries are bound in CHECKS.json.

The existing EP817 main was read at \nolinkurl{23c0110c95b5a2036bdc04a1f352b6e5e27742c8.} No existing main file, prior receipt, or accepted status is modified. The inspected Zeta interface is support\_diagrams.tex, equations (D6)--(D8), at \nolinkurl{58626a62cd5648e8ae7450fb54d4a4aab2330981,} blob \nolinkurl{d3493f891291ee6e94dbf2c77649f7d85d240df2.} The companion receiving note spells out the original diagrams and distinguishes word-value quotients from coefficient-amplitude quotients.

\section{References}\label{n14-references}

Frougny, C., \& Pelantová, E. (2018). \emph{Two applications of the spectrum of numbers} (Version 3; original preprint 2015) {[}Preprint{]}. arXiv:1512.04234v3. \url{https://arxiv.org/abs/1512.04234v3} . Retrieved abstract used for attribution of the prior finite-state Pisot methodology; no uninspected theorem body is an input to these proofs.

Korsky, S. (2026). \emph{Arithmetic progression-free subset-sum sets} (arXiv:2606.24139v1) {[}Preprint{]}. \url{https://arxiv.org/html/2606.24139v1} . Inspected Theorems 1.2--1.3 and Sections 4--6 for the general benchmark, the scalar ternary observation, and source-layout constraints.

KokunoYumeto. (2026). \emph{Reconstruction with changes of support index} {[}TeX source{]}. Zeta Function Research Reader, revision \nolinkurl{58626a62cd5648e8ae7450fb54d4a4aab2330981,} workbenches/splitzero-tandem/tex/support\_diagrams.tex, (D6)--(D8). Used for the original supported quotient and additional observation-kernel interface.

The Clankers. (2026, September 15). \emph{Supported zero actions, exact defect transport, and all-rank deformations} {[}Predecessor research contribution{]}. Supplied EP817\_SUPPORTED\_DEFECT\_20260915 archive. Used as the retained source and methodological predecessor, not as a new independent validation claim.
